% p3-09-computational

\section{P3 §9. Testable Computational Program}

9. Testable Computational Program

                Figure (fig-p3-ch9-program). Computational-programme triptych --- Algorithm suite /
                                       Reproducibility / Falsification criteria.

  9.1 Overview and Design Principles
  The framework makes specific predictions testable by finite computation. The design
  follows the pre-registration principle: all falsification criteria are stated before running
  the experiments, so they cannot be post-hoc adjusted. The computational program is
  built around three primary objects:

  1. The catalogue MDL vector M (Definition 7.3);
  2. Phase transition exponent estimates muT(xi) for each of the 42 catalogue entries;
  3. Pellis expansion coefficients of the sanctioned Fibonacci/Lucas seeds (regression
  tests).

  All computations use the reference Python package pellistrinity (Appendix B) at 50
  decimal digits of working precision.

  9.2 Algorithm Suite
  Algorithm 9.1 (Pellis expansion to depth D).
  Input: target x with N = 50 significant digits, depth D.
  Output: coefficient sequence (c0, …, cD), remainder RD.
  Implementation: greedy floor-subtraction via mpmath at 50 digits (Appendix B).

  Algorithm 9.2 (Trinity lattice search to complexity L).
  Input: target x, complexity budget L.
  Output: optimal label lambda\textasciicircum{}, approximant T\textasciicircum{}, residual R.
  Implementation: enumerate labels (n,k,p,m,q) with ell(lambda) $\leq$ L via nested loops;
  apply logarithmic pruning to reduce search from O(2L) to O(2L/2) using the identity log
  T = log n + k log 3 + p log$\varphi$ + mlogpi + q and nearest-neighbour search in log-space.

  Algorithm 9.3 (MDL score computation).
  Input: label lambda\textasciicircum{}*, residual coefficient sequence c res.
  Output: MDL score per Definition 7.1.
  Implementation: direct application of the bit-complexity formulas of Sections 2.6 and
  7.1.

  Algorithm 9.4 (Phase transition exponent estimation).
  Input: target x, maximum complexity budget Lmax.
  Output: estimated muT(x) with standard error.
  Implementation: for each L = 5, …, Lmax, compute PiL(x) and |RL(x)|. Perform linear
  regression of log|RL| against ell(lambda\textasciicircum{}*L); the slope times -1/log$\varphi$-1 is muT.

  Algorithm 9.5 (Catalogue MDL vector computation).
  Input: 42 catalogue entries {(xi, lambdai)}.
  Output: catalogue MDL vector M $\in$ R42.
  Implementation: apply Algorithms 9.1--9.3 to each entry; report mean, variance, and
  outlier structure.

  9.3 Reproducibility Protocol
  All algorithms are implemented in the reference Python package pellistrinity.
  Computation uses:

  - mpmath >= 1.3.0, sympy >= 1.12, numpy >= 1.24, scipy >= 1.10;
  - Working precision: 50 decimal digits throughout;
  - Sanctioned seeds only: {F17, F18, F19, F20, F21, L7, L8}; forbidden STROBE seeds
  {42, 43, 44, 45} must not appear as RNG initialisation values;
  - Physical constants from current NIST CODATA 2022 recommended values (
  - Mathematical constants from Wolfram MathWorld (

  9.4 Validation and Falsification Criteria
  A representation (PiL(x), c res) is validated if:

  1. Reconstruction agrees with the known value to at least 20 significant decimal digits;
  2. MDL score is strictly less than 2 D trivial (the description length of the naive decimal
  expansion at the same precision);
  3. Phase transition exponent estimate converges (standard error < 0.05) over L = 5, …,
  Lmax.

  Pre-registered falsification criteria (from the PhD program's Appendix B Falsification
  Ledger):

\subsection{References}

- Vasilev, D. \textit{Flos Aureus Monograph: Trinity Constants}. DOI \href{https://zenodo.org/doi/10.5281/zenodo.19227877}{10.5281/zenodo.19227877}.
- CODATA 2018 recommended values: NIST \href{https://physics.nist.gov/cuu/Constants/}{https://physics.nist.gov/cuu/Constants/}.
- Heyrovska, R. Golden-angle FSC publications --- Heyrovský Institute of Physical Chemistry.
